\documentclass[11pt]{article}
\usepackage[margin=1in]{geometry}
\usepackage[T1]{fontenc}
\usepackage{lmodern}
\usepackage{amsmath,amssymb,amsthm}
\usepackage{booktabs,array}
\usepackage{microtype,xcolor}
\usepackage[numbers,sort&compress]{natbib}
\usepackage[colorlinks=true,linkcolor=blue!50!black,citecolor=green!40!black,urlcolor=blue!50!black]{hyperref}
\usepackage{url}

\newtheorem{theorem}{Theorem}[section]
\newtheorem{proposition}[theorem]{Proposition}
\newtheorem{lemma}[theorem]{Lemma}
\newtheorem{corollary}[theorem]{Corollary}
\theoremstyle{definition}
\newtheorem{definition}[theorem]{Definition}
\theoremstyle{remark}
\newtheorem{remark}[theorem]{Remark}

\providecommand{\GL}{\mathrm{GL}}
\providecommand{\gl}{\mathfrak{gl}}
\providecommand{\tr}{\operatorname{tr}}
\providecommand{\cfg}[1]{{\small\texttt{#1}}}
\providecommand{\KL}{D_{\mathrm{KL}}}
\newcommand{\E}{\mathbb E}
\newcommand{\R}{\mathbb R}
\newcommand{\N}{\mathcal N}
\newcommand{\Qb}{\bar Q}
\newcommand{\curr}{\mathcal J}

\title{Dynamic Links: Latent Link Generators\\in the Variational Free Energy}
\author{Robert C. Dennis}
\date{Working note, October 3, 2026; revised October 4, 2026}

\begin{document}
\maketitle

\begin{abstract}
The exponential links $\Theta_{ij}=U_ie^{X_{ij}}U_j^{-1}$ of the gauge-covariant attention construction are promoted from parameters to latent random variables. The generators $X_{ij}$ live on the edges of a depth-by-token lattice, receive a normalized prior, and are inferred together with the token states and routing labels, so that their dynamics are inference rather than a postulated mechanics. We write the joint model, derive its variational free energy in closed form for mean-field recognition, and box the resulting functional. Its link sector is an identity-anchored Gaussian prior plus a curvature energy on the elementary loops of the lattice. We show that every continuous conjugation-invariant curvature measure on the non-compact group $\GL(K)$ is blind to unipotent holonomy, and we use instead the Frobenius defect of the reference-coordinate holonomy, which separates every holonomy from the identity and equals Wilson's plaquette density on $O(K)$. Stationarity gives a link field equation whose source is a precision-weighted prediction error times the source state, the delta rule of fast-weight learning. Without the link prior the optimal transport is rank one and lies outside the exponential chart; generators are then stationary only at critical points of the exponential map, and the infimum over the chart is either approached at that singular transport or lies strictly above it, so the prior is necessary. Linearized, the equation is a screened, massive discrete Maxwell equation of Euclidean signature, in which the time-like loops of the lattice supply the kinetic term and its coefficient, the precision of the implied depth increment, acts as an inertial mass for the maximum a posteriori path. Inference itself is first-order relaxation. Finally, directed (causal) matter induces no action on the links, while integrating out undirected Gaussian matter induces the action $\tfrac12\log\det J(X)$, bounded below, whose hopping expansion begins with Wilson's triangle action. All identities were checked in float64 at small size; nothing was trained.
\end{abstract}

\section{Scope}
\label{dl:sec:scope}

The exponential link $\Theta_{ij}=U_ie^{X_{ij}}U_j^{-1}$ is taken here as the primary construction. VFE 3.0 implements its static, shared-generator form, selected by \cfg{content\_link\_chart="exponential"}, in which the generator is a single learned matrix $X=OT$, the same on every edge. This note asks what changes when each $X_{ij}$ becomes a latent variable inferred from the context, so that the links evolve. The physics analogy is lattice gauge theory, in which link variables carry a Boltzmann weight $e^{-S}$ \cite{dl:wilson1974,dl:kogut1975,dl:creutz1983}. In the Euclidean setting that weight is a probability measure, so a variational free energy over latent links is the natural home for link dynamics; no separate Lagrangian is needed. The kinetic and curvature terms then appear as parts of $-\log$ of a declared prior over link fields, and the remaining question is which prior. Sections~\ref{dl:sec:curvature} and~\ref{dl:sec:induced} address that question from two sides: what a curvature energy on a non-compact group must look like, and what curvature energy matter induces when it is integrated out.

The scope is fixed throughout. The generative model is the layered causal model of Section~\ref{dl:sec:model}, with depth as the time axis. Recognition is mean field over token states, routing labels and links. The checks of Appendix~\ref{dl:app:checks} are float64 computations at $K\le4$ with a few tokens; they ground identities and code paths and do not set the direction of the theory, and no statement here is a training result.

\section{The depth-by-token lattice}
\label{dl:sec:setting}

Tokens are indexed by $i\in\{0,\dots,N-1\}$ with causal support $\mathcal S_i=\{j:j\le i\}$, and layers by $\ell\in\{0,\dots,L\}$. Each token carries a frame $U_i\in\GL(K)$, built once per forward pass and shared across layers as in the current architecture. The latent state of token $i$ at layer $\ell$ is written in reference coordinates, $\zeta_i^\ell\in\R^K$, with local coordinates $U_i\zeta_i^\ell$. A passive frame change $U_i\mapsto g_iU_i$ leaves every reference-coordinate quantity unchanged; this is the gauge convention used throughout.

A link edge connects the vertex $(j,\ell)$ to $(i,\ell+1)$ for every $j\in\mathcal S_i$. The edge carries a generator $X_{ij}^\ell\in\gl(K)$ in reference coordinates, its group element $G_{ij}^\ell=e^{X_{ij}^\ell}$, and the local link
\begin{equation}
\Theta_{ij}^\ell=U_iG_{ij}^\ell U_j^{-1},\qquad \Theta_{ij}^\ell\mapsto g_i\Theta_{ij}^\ell g_j^{-1}.
\label{dl:eq:link}
\end{equation}
The self edge $(i,\ell)\to(i,\ell+1)$ is time-like and the edges with $j<i$ are cross-token. Two-step routes from $(k,\ell)$ to $(i,\ell+2)$ pass through $(j,\ell+1)$ with $k\le j\le i$. We compare each route with the one that first stays at $k$, and define for $0\le\ell\le L-2$
\begin{equation}
W_{ijk}^\ell=\big(G_{ik}^{\ell+1}G_{kk}^\ell\big)^{-1}G_{ij}^{\ell+1}G_{jk}^\ell,\qquad k<j\le i.
\label{dl:eq:loops}
\end{equation}
The loop is trivial for $j=k$. For $j=i$ it is time-like: it compares moving at layer $\ell$ and staying at $\ell+1$ with staying first and moving second. For $k<j<i$ it is space-like. Any two routes through $j$ and $j'$ differ by $(W_{ij'k}^\ell)^{-1}W_{ijk}^\ell$, so the loops~\eqref{dl:eq:loops} generate all route comparisons between layers $\ell$ and $\ell+2$. Because $X$ is written in reference coordinates, $W_{ijk}^\ell$ itself, and not only its conjugacy class, is invariant under the passive action; the local holonomy is $U_kW_{ijk}^\ell U_k^{-1}$. When the links do not depend on the layer and the time-like links are trivial, $W_{ijk}^\ell$ reduces to the inverse-edge triangle holonomy $F_{ijk}=G_{ik}^{-1}G_{ij}G_{jk}$ of a layer-independent link field.

\section{Measuring curvature on a non-compact group}
\label{dl:sec:curvature}

A curvature energy must vanish on flat configurations and be positive otherwise. On compact groups the conjugation-invariant functions do this; on $\GL(K)$ they cannot.

\begin{proposition}[Class functions are blind to unipotent holonomy]
\label{dl:prop:blind}
Let $K\ge2$ and let $f:\GL(K,\R)\to\R$ be continuous and constant on conjugacy classes. Then $f(u)=f(I)$ for every unipotent $u$.
\end{proposition}

\begin{proof}
Write $u=RJR^{-1}$ with $J$ in Jordan form, so that $J=I+N$ with $N$ nilpotent and supported on the superdiagonal. Let $D_t=\operatorname{diag}(1,t^{-1},t^{-2},\dots,t^{-(K-1)})$. Then $(D_tJD_t^{-1})_{p,p+1}=tJ_{p,p+1}$, so $D_tJD_t^{-1}\to I$ as $t\to0$, and $f(u)=f(D_tR^{-1}uRD_t^{-1})\to f(I)$.
\end{proof}

The same scaling applied to the real Jordan form shows that such an $f$ depends only on the semisimple part of its argument, that is, on the eigenvalues. Hence the eigenvalue density $\sum_k|\lambda_k(W)-1|$, the spectral curvature $(\tfrac1K\sum_k|\mathrm{Log}\lambda_k|^2)^{1/2}$ of the implemented routing-holonomy diagnostic, and every other eigenvalue statistic vanish on the shear $W=I+E_{12}\ne I$ (Check~1). In particular, the bounded eigenvalue density of the multi-agent code does not vanish only at $W=I$.

\begin{definition}[Frobenius defect]
\label{dl:def:frob}
For a reference-coordinate holonomy $W$, set $\Phi(W)=\tfrac12\|W-I\|_F^2$.
\end{definition}

\begin{proposition}[Properties of the Frobenius defect]
\label{dl:prop:frob}
The defect is nonnegative and vanishes exactly at $W=I$. On $O(K)$ it equals Wilson's plaquette density, $\Phi(W)=K-\tr W$. It is invariant under the passive gauge action. Under the reparameterization $U_k\mapsto U_kh_k^{-1}$, $G_{ab}^\ell\mapsto h_aG_{ab}^\ell h_b^{-1}$, which leaves every local link unchanged, the holonomy transforms as $W_{ijk}^\ell\mapsto h_kW_{ijk}^\ell h_k^{-1}$, and $\Phi(h_kWh_k^{-1})=\Phi(W)$ for every $W\in\GL(K)$ exactly when $h_k$ is a nonzero multiple of an orthogonal matrix. For small generators,
\begin{equation}
\Phi(W_{ijk}^\ell)=\tfrac12\|(\delta X)_{ijk}^\ell\|_F^2+O(\|X\|^3),\qquad
(\delta X)_{ijk}^\ell=X_{ij}^{\ell+1}+X_{jk}^\ell-X_{ik}^{\ell+1}-X_{kk}^\ell.
\label{dl:eq:delta}
\end{equation}
For a time-like loop, $j=i$, this is $(\delta X)_{iik}^\ell=-\big[(X_{ik}^{\ell+1}-X_{ik}^\ell)-(X_{ii}^{\ell+1}-X_{kk}^\ell)\big]$, a discrete electric field: the depth difference of a cross-token link corrected by the time-like links at its ends.
\end{proposition}

\begin{proof}
Nonnegativity and the zero set are properties of the norm. For orthogonal $W$, $\|W-I\|_F^2=\tr(W^\top W)-2\tr W+K=2(K-\tr W)$. Passive invariance holds because $W$ is reference data. Under the reparameterization each two-step route from $(k,\ell)$ to $(i,\ell+2)$ becomes $h_i(\cdot)h_k^{-1}$, so $W\mapsto h_kWh_k^{-1}$. If $h_k=cO$ with $c\ne0$ and $O$ orthogonal, then $h_k(W-I)h_k^{-1}=O(W-I)O^\top$ has the Frobenius norm of $W-I$. Conversely, for $u,v\ne0$ and small $t\ne0$ the matrix $W=I+tuv^\top$ is invertible, and $\|h_k(W-I)h_k^{-1}\|_F=|t||h_ku||h_k^{-\top}v|$, so invariance gives $|h_ku|/|u|=|v|/|h_k^{-\top}v|$ for all nonzero $u$ and $v$. Both sides are then one constant $c>0$, so $|h_ku|=c|u|$ for every $u$ and $h_k/c$ is orthogonal. The expansion follows from $W=(I-X_{ik}^{\ell+1}-X_{kk}^\ell)(I+X_{ij}^{\ell+1}+X_{jk}^\ell)+O(\|X\|^2)$.
\end{proof}

The Frobenius defect thus uses one piece of structure that class functions do not: the reference metric. That metric reduces the measurement of curvature from $\GL(K)$ to its compact subgroup $O(K)$, where class functions would suffice: the reparameterizations that preserve $\Phi$ are the nonzero multiples of orthogonal matrices, and the scalar factor acts trivially on $W$. In local coordinates the metric is carried by the frame field. The same metric reappears below as the isotropic reference noise of the transport kernel and as the Frobenius anchor of the link prior; all three must be the same metric for the model to have the residual $O(K)$ symmetry of Proposition~\ref{dl:prop:frob}. A uniform scalar reparameterization leaves every link and the anchor unchanged but rescales the reference states, and the transport noise and the token priors fix that scale. The version of Proposition~\ref{dl:prop:frob} dated October 3, 2026 stated invariance exactly for orthogonal $h_k$, which omitted the scalar multiples.

\section{The generative model with latent links}
\label{dl:sec:model}

\begin{definition}[Latent-link model]
\label{dl:def:model}
The joint density of observations $o$, states $\zeta$, head labels $J$ and links $X$ is
\begin{equation}
p(o,\zeta,J,X)=p(X)\prod_ip(\zeta_i^0)\prod_{\ell=0}^{L-1}\prod_i\Big[\prod_h\pi_{i,J_i^{\ell,h}}^{\ell,h}\Big]\N\big(\zeta_i^{\ell+1};m_i^\ell,\Qb\big)\prod_ip(o_i\mid U_i\zeta_i^L),
\label{dl:eq:joint}
\end{equation}
with transport mean and link prior
\begin{equation}
m_i^\ell=\sum_hG_{i,J_i^{\ell,h}}^\ell P_h\zeta_{J_i^{\ell,h}}^\ell,\qquad
p(X)=Z_X^{-1}e^{-S_X(X)}.
\label{dl:eq:mean}
\end{equation}
Here $p(\zeta_i^0)=\N(\xi^p_{x_i},C^p_{x_i})$ is the reference-coordinate prior of token type $x_i$, $\pi^{\ell,h}$ is the causal routing prior, $\Qb\succ0$ is the reference transport noise, and $P_h=T^{-1}D_hT$ are the head projectors of the tiled heads, with $T$ the invertible stacked sending map and $D_h$ the coordinate projector onto the block of head $h$, so that $P_hP_g=\delta_{hg}P_h$ and $\sum_hP_h=I$. The link action is
\begin{equation}
S_X(X)=\sum_{\ell=0}^{L-1}\sum_i\sum_{j\in\mathcal S_i}\frac{\|X_{ij}^\ell-\bar X^\ell\|_F^2}{2\sigma^2}
+\lambda\sum_{\ell=0}^{L-2}\sum_{k<j\le i}w_{ijk}^\ell\Phi(W_{ijk}^\ell),\qquad
w_{ijk}^\ell=\bar\pi_{ij}^{\ell+1}\bar\pi_{jk}^{\ell},
\label{dl:eq:action}
\end{equation}
with $\bar\pi$ the head-averaged routing prior, an anchor $\bar X^\ell$, anchor width $\sigma>0$ and curvature coupling $\lambda\ge0$.
\end{definition}

Each head carries the component $P_h\zeta_j$ that it reads from its own source and transports it through that source's link, so the link acts after the head reads. The path weights in~\eqref{dl:eq:action} use the generative routing prior and never the recognition weights, in keeping with the modeling rule, adopted here, that no generative factor reads a recognition law. The model is properly normalized: $S_X\ge\|X-\bar X\|_F^2/(2\sigma^2)$, so $Z_X$ is finite, and every other factor in~\eqref{dl:eq:joint} is a normalized conditional. All link-sector quantities are reference data, so the passive gauge action changes only the frames, which enter the joint only through the decoder.

\begin{remark}[Static limit]
\label{dl:rem:static}
If $\bar X^\ell=OT$ for every layer and $\sigma\to0$, every link is pinned to the shared generator and the mean of $m_i^\ell$ under routing weights $\beta$ is $e^{OT}T^{-1}B_{ij}T\xi_j$ summed over $j$, with the head-diagonal routing matrix $B_{ij}=\sum_h\beta_{ij}^{\ell,h}D_h$, the exponential-chart message with shared generator $OT$ (Check~10). The implemented \cfg{content\_link\_chart="exponential"} is therefore the zero-temperature limit of this model, and the dynamic links are the fluctuations around it. Per-head links $X_{ij}^{\ell,h}$ are a variant that this note does not pursue.
\end{remark}

\section{The variational free energy}
\label{dl:sec:vfe}

Recognition is mean field,
\begin{equation}
q(\zeta,J,X)=\prod_{\ell,i}\N(\zeta_i^\ell;\xi_i^\ell,C_i^\ell)\prod_{\ell,i,h}\beta_i^{\ell,h}(J_i^{\ell,h})q(X),
\label{dl:eq:recognition}
\end{equation}
with $q(X)$ either a point mass at $M$, which gives maximum a posteriori links, or a Gaussian on $\gl(K)^{\text{edges}}$. With $A=\Qb^{-1}$, write the routed moments of the transport mean under $q$,
\begin{equation}
\bar m_i^\ell=\E_q[m_i^\ell],\qquad V_i^\ell=\operatorname{Cov}_q(m_i^\ell).
\label{dl:eq:moments}
\end{equation}
For a point mass $q(X)$, with $a_{ij}^h=G_{ij}^\ell P_h\xi_j^\ell$, $\bar m_i^h=\sum_j\beta_{ij}^{\ell,h}a_{ij}^h$, $c^{hh}_{ij}=\beta^{\ell,h}_{ij}$ and $c^{hg}_{ij}=\beta^{\ell,h}_{ij}\beta^{\ell,g}_{ij}$ for $h\ne g$,
\begin{equation}
V_i^\ell=\sum_h\sum_j\beta_{ij}^{\ell,h}(a_{ij}^h-\bar m_i^h)(a_{ij}^h-\bar m_i^h)^\top+\sum_{h,g}\sum_jc^{hg}_{ij}G_{ij}^\ell P_hC_j^\ell P_g^\top G_{ij}^{\ell\top},
\label{dl:eq:V}
\end{equation}
the moment-matched predictive covariance of the current implementation without its floor.

\begin{lemma}[Transport term]
\label{dl:lem:transport}
Under~\eqref{dl:eq:recognition},
\begin{equation}
\E_q\big[-\log\N(\zeta_i^{\ell+1};m_i^\ell,\Qb)\big]-H\big[q(\zeta_i^{\ell+1})\big]
=\KL\big(q(\zeta_i^{\ell+1})\|\N(\bar m_i^\ell,\Qb)\big)+\tfrac12\tr\big(AV_i^\ell\big).
\label{dl:eq:transport}
\end{equation}
For one head and a point mass $q(X)$ the right side equals
\begin{equation}
\sum_j\beta_{ij}^\ell\Big[\KL\big(q(\zeta_i^{\ell+1})\|\N(G_{ij}^\ell\xi_j^\ell,\Qb)\big)+\tfrac12\tr\big(AG_{ij}^\ell C_j^\ell G_{ij}^{\ell\top}\big)\Big].
\label{dl:eq:transport_onehead}
\end{equation}
\end{lemma}

\begin{proof}
Under mean field $\zeta_i^{\ell+1}$ is independent of $m_i^\ell$, so $\E_q[(\zeta-m)(\zeta-m)^\top]=C_i^{\ell+1}+(\xi_i^{\ell+1}-\bar m_i^\ell)(\xi_i^{\ell+1}-\bar m_i^\ell)^\top+V_i^\ell$. Inserting this into the Gaussian cross entropy and subtracting the entropy $\tfrac12\log\det(2\pi eC_i^{\ell+1})$ gives~\eqref{dl:eq:transport}. For one head, $V_i^\ell=\sum_j\beta_{ij}G_{ij}C_jG_{ij}^\top+\sum_j\beta_{ij}(a_j-\bar m)(a_j-\bar m)^\top$, and for Gaussians with a common covariance $\sum_j\beta_j\KL(q\|\N(a_j,\Qb))=\KL(q\|\N(\bar m,\Qb))+\tfrac12\sum_j\beta_j\|a_j-\bar m\|_A^2$, which gives~\eqref{dl:eq:transport_onehead}.
\end{proof}

Equation~\eqref{dl:eq:transport_onehead} identifies the matter coupling: a routing-weighted alignment between the receiver's belief and each transported source, in which the source's uncertainty is pushed through the link and charged in noise units. It is the expected-likelihood (channel) form of the alignment energy rather than the ambient sum of divergences, because the joint~\eqref{dl:eq:joint} is a likelihood model. Collecting the factors of~\eqref{dl:eq:joint} gives the free energy.

\begin{equation}
\boxed{
\begin{aligned}
\mathcal F[q]={}&\sum_i\E_q\big[-\log p(o_i\mid U_i\zeta_i^L)\big]+\sum_i\KL\big(q(\zeta_i^0)\|p(\zeta_i^0)\big)\\
&+\sum_{\ell=0}^{L-1}\sum_i\Big[\KL\big(q(\zeta_i^{\ell+1})\|\N(\bar m_i^\ell,\Qb)\big)+\tfrac12\tr\big(\Qb^{-1}V_i^\ell\big)\Big]
+\sum_{\ell,i,h}\KL\big(\beta_i^{\ell,h}\|\pi_i^{\ell,h}\big)\\
&+\E_{q(X)}\Big[\sum_{\ell,i,j\in\mathcal S_i}\frac{\|X_{ij}^\ell-\bar X^\ell\|_F^2}{2\sigma^2}
+\lambda\sum_{\ell}\sum_{k<j\le i}\bar\pi_{ij}^{\ell+1}\bar\pi_{jk}^\ell\tfrac12\big\|W_{ijk}^\ell-I\big\|_F^2\Big]-H\big[q(X)\big]+\log Z_X
\end{aligned}}
\label{dl:eq:F}
\end{equation}

The first line is the evidence and encoding sector, the second the matter sector of Lemma~\ref{dl:lem:transport} together with the routing sector, and the third the link sector, $\KL(q(X)\|p(X))$, with the anchor acting as a mass term and the loop energy as the curvature term. By Jensen's inequality $\mathcal F[q]\ge-\log p(o)$ for every $q$, so~\eqref{dl:eq:F} is an exact variational bound for the declared model. For a point mass $q(X)=\delta_M$ the entropy and $\log Z_X$ are constants and the third line is $S_X(M)$. Because the links enter $\bar m$ and $V$, every sector of the second line couples to them.

\section{The link field equation}
\label{dl:sec:field}

Write $L(X,E)=\int_0^1e^{sX}Ee^{(1-s)X}\mathrm{d}s$ for the Fr\'echet derivative of the matrix exponential \cite{dl:higham2008}. Its Frobenius adjoint is $L(X^\top,\cdot)$, so $\partial f(e^X)/\partial X=L(X^\top,\partial f/\partial G)$.

\begin{proposition}[Field equation]
\label{dl:prop:field}
For point-mass links, $\mathcal F$ is stationary in $X_{ij}^\ell$ exactly when
\begin{equation}
\frac{\partial S_X}{\partial X_{ij}^\ell}=L\big(X_{ij}^{\ell\top},\curr_{ij}^\ell\big),\qquad
\curr_{ij}^\ell=\sum_h\beta_{ij}^{\ell,h}A\E_q\big[(\zeta_i^{\ell+1}-m_i^\ell)\zeta_j^{\ell\top}\mid J_i^{\ell,h}=j\big]P_h^\top.
\label{dl:eq:field}
\end{equation}
For one head the matter current is
\begin{equation}
\curr_{ij}^\ell=\beta_{ij}^\ell A\Big[\big(\xi_i^{\ell+1}-G_{ij}^\ell\xi_j^\ell\big)\xi_j^{\ell\top}-G_{ij}^\ell C_j^\ell\Big].
\label{dl:eq:current}
\end{equation}
\end{proposition}

\begin{proof}
Only the transport term of receiver $i$ at layer $\ell$ and the link action depend on $X_{ij}^\ell$. By Lemma~\ref{dl:lem:transport} the transport term is $\tfrac12\E_q\|\zeta_i^{\ell+1}-m_i^\ell\|_A^2$ up to $X$-free terms. Since $m_i^\ell$ contains $G_{ij}^\ell P_h\zeta_j^\ell$ exactly on the event $J_i^{\ell,h}=j$, its derivative in $G_{ij}^\ell$ is $-\curr_{ij}^\ell$, and the chain rule through $e^X$ gives~\eqref{dl:eq:field}. For one head, $\E[(\zeta'-m)\zeta_j^\top\mid J=j]=\xi'\xi_j^\top-G_{ij}(C_j+\xi_j\xi_j^\top)$, which is~\eqref{dl:eq:current}.
\end{proof}

The current is a precision-weighted prediction error times the source state, minus a term from the source's uncertainty. It is the delta rule, the update of fast-weight memories \cite{dl:schmidhuber1992}, here derived from a free energy and filtered through the derivative of the exponential chart. The left side of~\eqref{dl:eq:field} is the response of the link action. The equation has the form of a field equation with a source: link stress equals matter current.

\begin{corollary}[The link prior is necessary]
\label{dl:cor:prior}
Let $S_X=0$, one head, $K\ge2$ and $\beta_{ij}^\ell>0$, and write $\xi'=\xi_i^{\ell+1}$, $S_j^\ell=C_j^\ell+\xi_j^\ell\xi_j^{\ell\top}$ and $\mathcal T_{ij}(G)=\tfrac12\beta_{ij}^\ell\E_q\|\zeta_i^{\ell+1}-G\zeta_j^\ell\|_A^2$ for real $K\times K$ matrices $G$. Up to an $X$-free constant, $\mathcal F$ depends on $X_{ij}^\ell$ through $\mathcal T_{ij}(e^{X_{ij}^\ell})$, and:
\begin{enumerate}
\item[(i)] $\mathcal T_{ij}$ is a strictly convex quadratic whose only stationary point is $G^\star=\xi'\xi_j^{\ell\top}(S_j^\ell)^{-1}$, of rank at most one and therefore outside the exponential chart, since $\det e^X=e^{\tr X}>0$.
\item[(ii)] A generator $X_{ij}^\ell$ is stationary only where $L(X_{ij}^{\ell\top},\cdot)$ is singular, at a critical point of the exponential map. Such generators exist: for $K=2$, $\xi'=\xi_j^\ell=e_1$ and $C_j^\ell=A=I$, the generator $X=-(\log3)I+2\pi\bigl(\begin{smallmatrix}0&-1\\1&0\end{smallmatrix}\bigr)$ is stationary with link $\tfrac13I\ne G^\star$.
\item[(iii)] If $\xi_j^{\ell\top}(C_j^\ell)^{-1}\xi'\ge0$, the infimum of $\mathcal T_{ij}$ over the exponential chart equals $\mathcal T_{ij}(G^\star)$, is not attained, and is approached only along links converging to $G^\star$.
\item[(iv)] If $\xi_j^{\ell\top}(C_j^\ell)^{-1}\xi'<0$, no sequence of chart links converges to $G^\star$, and the infimum over the chart lies strictly above $\mathcal T_{ij}(G^\star)$. For $K=2$ and $K=3$ it is not attained.
\end{enumerate}
\end{corollary}

\begin{proof}
By Lemma~\ref{dl:lem:transport} the $X_{ij}^\ell$-dependent part of $\mathcal F$ is $\mathcal T_{ij}(e^{X_{ij}^\ell})$ up to a constant. The gradient of $\mathcal T_{ij}$ is $\beta_{ij}^\ell A(GS_j^\ell-\xi'\xi_j^{\ell\top})$ and its Hessian $E\mapsto\beta_{ij}^\ell AES_j^\ell$ is positive definite, as in the proof of Proposition~\ref{dl:prop:linear}; indeed $\mathcal T_{ij}(G)-\mathcal T_{ij}(G^\star)=\tfrac12\beta_{ij}^\ell\|A^{1/2}(G-G^\star)(S_j^\ell)^{1/2}\|_F^2$. This gives (i). By Proposition~\ref{dl:prop:field}, $X$ is stationary exactly when $L(X^\top,\nabla\mathcal T_{ij}(e^X))=0$, and $\nabla\mathcal T_{ij}(e^X)\ne0$ because $e^X\ne G^\star$, so $L(X^\top,\cdot)$ is singular. In the example, $\nabla\mathcal T_{ij}(\tfrac13I)=\tfrac13\beta_{ij}^\ell\operatorname{diag}(-1,1)$, and $\operatorname{diag}(-1,1)$ anticommutes with the rotation generator $J$, and $L(X^\top,\operatorname{diag}(-1,1))=\tfrac13\operatorname{diag}(-1,1)\int_0^1e^{2\pi(1-2s)J}\mathrm ds=0$. For (iii) and (iv), write $\tau=\tr G^\star=\xi_j^{\ell\top}(S_j^\ell)^{-1}\xi'$, which by the Sherman-Morrison formula is $\xi_j^{\ell\top}(C_j^\ell)^{-1}\xi'$ divided by $1+\xi_j^{\ell\top}(C_j^\ell)^{-1}\xi_j^\ell$. The eigenvalues of $G^\star$ are $\tau$ and $0$ with multiplicity $K-1$. A real matrix is the exponential of a real matrix exactly when it is invertible and every Jordan block belonging to a negative eigenvalue occurs an even number of times \cite[Theorem~1]{dl:culver1966}. If $\tau\ge0$, then $G^\star+\varepsilon I$ has only positive eigenvalues for $\varepsilon>0$, so it is a chart link converging to $G^\star$, and by the quadratic identity every sequence of links whose energy tends to $\mathcal T_{ij}(G^\star)$ converges to $G^\star$. If $\tau<0$, it is a simple eigenvalue of $G^\star$, so every real matrix near $G^\star$ has a simple negative eigenvalue near $\tau$, a single Jordan block of size one; such a matrix is no chart link. A neighborhood of $G^\star$ therefore contains no chart link, and outside it the quadratic identity keeps $\mathcal T_{ij}$ above $\mathcal T_{ij}(G^\star)$ by a positive amount. Non-attainment for $K=2,3$ follows from Proposition~\ref{dl:prop:semisimple} of Appendix~\ref{dl:app:proofs}, applied to $\mathcal T_{ij}$, whose only stationary point $G^\star$ is singular; its perturbation argument covers every chart link at $K\le3$.
\end{proof}

The version of this corollary dated October 3, 2026 stated that the infimum is always approached at the singular boundary. That holds only in case (iii). In case (iv), which a single context reaches on an open set of states, the chart infimum lies at strictly higher energy, and whether it is attained for $K\ge4$, at a chart link with a negative eigenvalue that has a Jordan block of size at least two, the only chart links that Proposition~\ref{dl:prop:semisimple} does not exclude, is open. In every case the stationary generators of the unregularized problem, where they exist, sit at critical points of the exponential map.

A single context supplies one source state per edge, which determines only a rank-one regression. Averaged over a population of contexts, the same normal equation gives, for a centered population law, the population regression link $C_{ab}C_{bb}^{-1}$, with $C_{ab}$ the population cross-covariance of the states at the vertices $a$ and $b$, whose curvature measures conditional dependence (Proposition~\ref{dl:prop:regression} of Appendix~\ref{dl:app:proofs}). The anchor $\sigma$ is what makes the per-context problem well posed, and its mean $\bar X$ is the learned, slow part of the link.

\begin{proposition}[Linearized field equation]
\label{dl:prop:linear}
Take one head, a flat anchor $\bar X^\ell=0$, small links and small sources. To first order the field equation is
\begin{equation}
\frac1{\sigma^2}X_e+\lambda(\delta^\top w\delta X)_e+\beta_eAX_eS_{j(e)}
=\beta_eA\big[(\xi_i^{\ell+1}-\xi_j^\ell)\xi_j^{\ell\top}-C_j^\ell\big],
\label{dl:eq:linear}
\end{equation}
for every edge $e=(i\leftarrow j)$ at layer $\ell$, where $\delta$ is the loop map~\eqref{dl:eq:delta} and $w$ the diagonal of loop weights. The operator on the left is self-adjoint in the Frobenius inner product and bounded below by $\sigma^{-2}$, so the linearized links are unique.
\end{proposition}

\begin{proof}
At $X=0$, $L(0,\cdot)$ is the identity and $G=I+X+O(X^2)$, so~\eqref{dl:eq:current} is $\beta A[(\xi'-\xi_j)\xi_j^\top-C_j-XS_j]$ plus terms of second order. By~\eqref{dl:eq:delta} the loop energy is $\tfrac\lambda2\sum w\|\delta X\|_F^2+O(X^3)$, whose gradient is $\lambda\delta^\top w\delta X$. The map $X\mapsto AXS$ is self-adjoint and nonnegative because $\langle Y,AXS\rangle=\tr(Y^\top AXS)=\langle AYS,X\rangle$ and $\langle X,AXS\rangle=\|A^{1/2}XS^{1/2}\|_F^2$, and $\delta^\top w\delta$ is nonnegative.
\end{proof}

Read as physics, \eqref{dl:eq:linear} is a static massive Maxwell (Proca) equation on the depth-by-token complex. The anchor is a bare mass, the loop term is the curl-curl operator, and the matter term $\beta AXS$ is an additional, state-dependent mass, analogous to the screening a charged medium gives the photon. The source is the prediction error times the source state. Through the time-like loops, the loop term contains $\tfrac\lambda2w_{iik}^\ell\|(X_{ik}^{\ell+1}-X_{ik}^\ell)-(X_{ii}^{\ell+1}-X_{kk}^\ell)\|_F^2$, the square of a discrete electric field. The kinetic term is therefore not a separate postulate: it is the time-like part of the same curvature energy, exactly as the electric term arises from time-like plaquettes in lattice gauge theory \cite{dl:kogut1975}, and its coefficient $\lambda\bar\pi_{ii}^{\ell+1}\bar\pi_{ik}^\ell$ is the curvature coupling times the prior probability of the route. The space-like loops give the magnetic part. The operator is positive, so the equation is elliptic: depth enters with the same sign as the token directions, the Euclidean signature (Check~6 compares~\eqref{dl:eq:linear} with the exact stationary links).

\section{Dynamics}
\label{dl:sec:dynamics}

Two notions of dynamics follow from~\eqref{dl:eq:F}, and they should not be confused.

\paragraph{Inference dynamics.} Within a forward pass the links relax toward the stationary point of $\mathcal F$. For point-mass links the natural preconditioner is the Fisher information of the transported law with respect to the generator, which for one head is
\begin{equation}
\mathcal I_{ij}^\ell[E,E']=\beta_{ij}^\ell\tr\big(AL(X_{ij}^\ell,E)S_j^\ell L(X_{ij}^\ell,E')^\top\big),
\label{dl:eq:fisher}
\end{equation}
the Gauss-Newton matrix of the transport term; because the kernel covariance $\Qb$ does not depend on the link, only the mean part of the Gaussian Fisher metric contributes. The flow $\dot M=-\mathcal I(M)^{-1}\nabla_M\mathcal F$ is first order, with the Fisher information as friction and no inertia (Lemma~\ref{dl:lem:fisher} of Appendix~\ref{dl:app:proofs}). For Gaussian $q(X)=\N(M,\Sigma_X)$ the corresponding natural-gradient updates are those of the Bayesian learning rule \cite{dl:khan2023}, $\Sigma_X^{-1}\leftarrow(1-\rho)\Sigma_X^{-1}+\rho\E_q[\nabla^2\ell]$ and $M\leftarrow M-\rho\Sigma_X\E_q[\nabla\ell]$, where $\ell$ is the $X$-dependent part of $\mathcal F$ before the entropy. In the Laplace approximation $\Sigma_X^{-1}=\mathcal I+\nabla^2S_X$ at the stationary point, and the link entropy contributes an Occam term to~\eqref{dl:eq:F}.

\paragraph{Depth as time, and where inertia comes from.} The maximum a posteriori link field over all layers solves~\eqref{dl:eq:field} jointly in depth. The following statement makes the intuition of precision as inertial mass exact in the one place where it holds.

\begin{proposition}[Inertia of the maximum a posteriori path]
\label{dl:prop:inertia}
Let a path $X^0,\dots,X^L$ with fixed ends minimize $\sum_\ell\big[\tfrac m2\|X^{\ell+1}-X^\ell\|_F^2+V(X^\ell)\big]$ with $V$ differentiable. Then
\begin{equation}
m\big(2X^\ell-X^{\ell+1}-X^{\ell-1}\big)+\nabla V(X^\ell)=0,
\label{dl:eq:inertia}
\end{equation}
the discretization of $m\ddot X=\nabla V$: Newton's equation with mass $m$ in the potential $-V$.
\end{proposition}

\begin{proof}
Differentiate the objective in $X^\ell$ for $0<\ell<L$.
\end{proof}

The quadratic increment term is $-\log$ of a Gaussian increment density with precision $m$, and in~\eqref{dl:eq:F} it is supplied by the time-like loops with $m=\lambda\bar\pi_{ii}^{\ell+1}\bar\pi_{ik}^\ell$ when the time-like links are held at the identity. Precision is therefore the inertial mass of the smoothed path. Two qualifications follow from the sign. The equation is the Euclidean one, motion in the inverted potential. For a scalar link coordinate held at the minima of a double well at its two ends, every maximum a posteriori path is nondecreasing in depth, and strictly increasing when the minima are critical points of the potential, so it is a monotone kink that does not oscillate (Proposition~\ref{dl:prop:kink} of Appendix~\ref{dl:app:proofs}; Check~9 illustrates it for $V(x)=(x^2-1)^2/4$). The argument uses the order of the real line and covers a one-dimensional link coordinate only. And the inertia belongs to the smoother, the joint optimization over depth; filtering layer by layer gives first-order dynamics. Real-time propagation, waves of the link field, would require an analytic continuation that no inference procedure supplies. Second-order dynamics in inference time can be obtained by declaring velocity states in the generative model, as in generalized coordinates of motion \cite{dl:friston2008dem}, which again ties the mass to a declared precision; the path-integral form of such increment priors is the Onsager-Machlup functional \cite{dl:onsager1953}.

\section{Where the curvature energy comes from}
\label{dl:sec:induced}

The loop energy in~\eqref{dl:eq:action} is a declared prior. Integrating out the token states shows when such an energy is induced instead.

\begin{proposition}[Directed matter induces no link action]
\label{dl:prop:directed}
In the model~\eqref{dl:eq:joint}, $\int p(\zeta,J\mid X)\mathrm{d}\zeta\mathrm{d}J=1$ for every $X$, so the prior marginal of the links is $p(X)$ itself.
\end{proposition}

\begin{proof}
Every factor other than $p(X)$ and the decoder is a normalized conditional of a directed acyclic model, and the decoder integrates to one over $o$.
\end{proof}

In the causal model, curvature reaches the posterior only through evidence. The predictive law of $\zeta_i^{\ell+2}$ given $\zeta_k^\ell$ is a mixture over routes $j$ of Gaussians centered at $G_{ik}^{\ell+1}G_{kk}^\ell W_{ijk}^\ell\zeta_k^\ell$, whose components coincide exactly when the loops are trivial. Evidence generated by a single transport therefore favors flat links, while evidence that requires different routes to carry different content favors curvature.

An induced action appears for undirected matter. Let the states of one layer interact through pairwise consistency factors,
\begin{equation}
p(\zeta\mid X)\propto\exp\Big(-\tfrac12\sum_v\zeta_v^\top\Lambda_0\zeta_v-\tfrac12\sum_{(i\leftarrow j)}\|\zeta_i-G_{ij}\zeta_j\|_\Lambda^2\Big),
\label{dl:eq:undirected}
\end{equation}
with $\Lambda_0,\Lambda\succ0$, the energy-based reading of the alignment terms. Its precision matrix $J(X)=D(X)-\mathcal A(X)$ has diagonal blocks $D_v=\Lambda_0+\sum_{v\leftarrow\cdot}\Lambda+\sum_{\cdot\leftarrow v}G^\top\Lambda G$ and off-diagonal blocks $\mathcal A_{ij}=\Lambda G_{ij}$, $\mathcal A_{ji}=G_{ij}^\top\Lambda$.

\begin{proposition}[Induced action of undirected Gaussian matter]
\label{dl:prop:induced}
Let the joint density of states and links be proportional to $\tilde p(X)e^{-E(\zeta,X)}$, with $E$ the exponent of~\eqref{dl:eq:undirected} and $\tilde p$ a base prior, normalized as a whole and not conditionally in $\zeta$. Then the link marginal is proportional to $\tilde p(X)e^{-S_{\mathrm{ind}}(X)}$, where
\begin{equation}
S_{\mathrm{ind}}(X)=\tfrac12\log\det J(X)\ge\tfrac12\log\det(I\otimes\Lambda_0).
\label{dl:eq:induced}
\end{equation}
If the spectral radius of $D^{-1}\mathcal A$ is below one,
\begin{equation}
S_{\mathrm{ind}}(X)=\tfrac12\log\det D-\tfrac12\sum_{n\ge1}\frac1n\tr\big[(D^{-1}\mathcal A)^n\big],
\label{dl:eq:hopping}
\end{equation}
and the $n$-th term is a sum over closed walks of length $n$ of traces of products of the blocks of $D^{-1}\mathcal A$ along the walk. Each block of $\mathcal A$ transports between adjacent sites, by $\Lambda G_{ij}$ when it follows the edge $(i\leftarrow j)$ from $j$ to $i$ and by the transpose $G_{ij}^\top\Lambda$ when it traverses that edge against its direction; for orthogonal links the transpose is the inverse, and these traces are Wilson loops. For isotropic $\Lambda_0=\lambda_0I$, $\Lambda=\lambda_EI$ and orthogonal links, $D_v=d_vI$ does not depend on the links, the walks of length two contribute constants, and on a triangle
\begin{equation}
S_{\mathrm{ind}}=\text{const}+\frac{\lambda_E^3}{d_id_jd_k}\big(K-\tr F_{ijk}\big)+O\Big(\frac{\lambda_E^5}{d^5}\Big),
\label{dl:eq:wilson}
\end{equation}
which is $\lambda_E^3/(d_id_jd_k)$ times the Frobenius defect $\Phi(F_{ijk})$.
\end{proposition}

\begin{proof}
Integrating $e^{-E(\zeta,X)}$ over $\zeta$ gives $(2\pi)^{nK/2}\det J(X)^{-1/2}$, so the link marginal is proportional to $\tilde p(X)\det J(X)^{-1/2}$. Had the states been given a conditionally normalized law, this factor would cancel, which is the content of Proposition~\ref{dl:prop:directed}. Each pairwise term is positive semidefinite, so $J\succeq I\otimes\Lambda_0$. The expansion is $\log\det(D-\mathcal A)=\log\det D+\log\det(I-D^{-1}\mathcal A)$ with the logarithmic series. Expanding the block product $\tr[(D^{-1}\mathcal A)^n]$ gives a sum over sequences of sites, and $\mathcal A_{vu}=0$ unless $v$ and $u$ are joined by an edge, with $\mathcal A_{ij}=\Lambda G_{ij}$ and $\mathcal A_{ji}=G_{ij}^\top\Lambda$ for the edge $(i\leftarrow j)$, which leaves the closed walks with the stated factors. With orthogonal links, $G^\top\Lambda G=\lambda_EI$, so $D$ is link independent, and a walk $i\to j\to i$ contributes $\lambda_E^2\tr(GG^\top)/(d_id_j)$, a constant. A closed walk around the triangle contributes $\lambda_E^3\tr F/(d_id_jd_k)$ in one orientation and $\lambda_E^3\tr F^{-1}=\lambda_E^3\tr F$ in the other; with three starting points and two orientations, the $n=3$ term is $-\tfrac12\cdot\tfrac13\cdot6\lambda_E^3\tr F/(d_id_jd_k)$. Proposition~\ref{dl:prop:frob} identifies $K-\tr F$ with $\Phi(F)$.
\end{proof}

This is the Gaussian-matter version of the hopping-parameter expansion of lattice gauge theory \cite{dl:montvay1994} and of induced gauge actions \cite{dl:kazakov1993}: integrating out matter generates the Wilson action at leading order, and the coupling is the cube of the hopping parameter. Check~8 confirms~\eqref{dl:eq:wilson}; the ratio of the exact induced action to the leading term is $1.15$, $1.022$, $1.003$ and $1.0003$ at $\lambda_0=3$, $10$, $30$ and $100$. The induced action also settles the non-compact directions. With orthogonal links on a tree it is constant, but a symmetric generator on a single edge changes it through the diagonal blocks even without loops, so the non-compact directions acquire edge-local mass terms while the compact directions enter only through loops. The induced action is bounded below by~\eqref{dl:eq:induced} for every $\GL(K)$ configuration, and it responds to a unipotent holonomy that every class function misses (Check~8). The causal predictive route is directed, so Proposition~\ref{dl:prop:directed} applies to it, and the loop energy of~\eqref{dl:eq:action} is a declared prior whose form, the Frobenius defect with its Wilson limit, is the one undirected Gaussian matter would induce.

\section{Relation to the implementation and next steps}
\label{dl:sec:next}

The implemented exponential chart is the static limit of Remark~\ref{dl:rem:static}. A first implementation of dynamic links would keep $\bar X=OT$ as the slow, learned anchor and infer per-context deviations $M-\bar X$ with a few Gauss-Newton steps preconditioned by~\eqref{dl:eq:fisher}, or with one solve of the linearized equation~\eqref{dl:eq:linear}. The cost per step is one exponential per edge, $O(LN^2K^3)$ over all causal edges, which calls for restricted supports, low-rank deviations or amortization, in which a network predicts $M-\bar X$ from the states; amortized deviations are a fast-weight program with the delta-rule source of Proposition~\ref{dl:prop:field}. A grounding experiment at small size would check that inferred links reproduce the stationary equation, that the loop energy responds to unipotent holonomy, and that removing the anchor drives links toward the singular transport when $\xi_j^{\ell\top}(C_j^\ell)^{-1}\xi_i^{\ell+1}\ge0$ and leaves them at strictly higher energy otherwise, as Corollary~\ref{dl:cor:prior} predicts. Such checks validate the construction; they do not decide whether dynamic links help a language model at scale.

Several questions remain open. The values of $\lambda$, $\sigma$ and the loop weights are modeling choices without a derivation here. The multi-head field equation couples edges through the receiver's mixture, and its linearization has cross-edge terms that this note does not analyze. Stationarity in the time-like links gives a constraint analogous to Gauss's law that remains to be worked out. Full Gaussian link posteriors cost $K^2\times K^2$ per edge. Whether evidence-driven curvature improves prediction, and whether the undirected matter model should replace the declared loop prior, are empirical and modeling questions respectively.

\appendix
\section{Supplementary proofs}
\label{dl:app:proofs}

This appendix proves four statements used in the main text: the Fisher form~\eqref{dl:eq:fisher} and the first-order character of the recognition flow, the non-attainment part of Corollary~\ref{dl:cor:prior}(iv), the curvature of population regression links, and the monotonicity of the maximum a posteriori path of a scalar link coordinate. Matrices carry the Frobenius inner product $\langle E,E'\rangle=\tr(E^\top E')$ and its norm, $S_j^\ell=C_j^\ell+\xi_j^\ell\xi_j^{\ell\top}$ is the second moment of $\zeta_j^\ell$ under the recognition law~\eqref{dl:eq:recognition}, and a chart link is a matrix $e^X$ with $X\in\gl(K)$. By \cite[Theorem~1]{dl:culver1966}, a real matrix is a chart link exactly when it is invertible and every Jordan block belonging to a negative eigenvalue occurs an even number of times; in particular, every invertible real matrix without negative eigenvalues is a chart link.

\begin{lemma}[Fisher information of the transported law]
\label{dl:lem:fisher}
Take one head and point-mass links. For an edge $(i\leftarrow j)$ at layer $\ell$ with generator $X=X_{ij}^\ell$, the Fisher information with respect to $X$ of the transport kernel $\N(\zeta';e^X\zeta_j^\ell,\Qb)$, averaged over $q(\zeta_j^\ell)$ and weighted by the routing probability $\beta_{ij}^\ell$, is the form~\eqref{dl:eq:fisher}. It equals the Gauss-Newton form of the transport term $\mathcal T_{ij}(e^X)$ of Corollary~\ref{dl:cor:prior}, only the mean part of the Gaussian Fisher metric contributes to it, and it is positive definite exactly when $\beta_{ij}^\ell>0$ and $L(X,\cdot)$ is invertible. On the open set where the block-diagonal form $\mathcal I(M)$ with these blocks is invertible, the recognition flow $\dot M=-\mathcal I(M)^{-1}\nabla_M\mathcal F$ of the point-mass links is of first order: its velocity is a function of the configuration alone and vanishes exactly at the critical points of $\mathcal F$, it obeys the force balance $\mathcal I(M)[\dot M,\cdot]=-d\mathcal F$ with $\mathcal I(M)$ as friction tensor and no inertial term, $\mathcal F$ strictly decreases along every nonconstant solution, so that no nonconstant solution is periodic, and at a critical point with invertible Hessian $H$ the linearized field $-\mathcal I^{-1}H$ has real eigenvalues.
\end{lemma}

\begin{proof}
The mean of the kernel is $e^X\zeta_j^\ell$, with derivative $L(X,E)\zeta_j^\ell$ in the direction $E$. For a Gaussian kernel with fixed covariance $\Qb$, the score in the direction $E$ is $\langle L(X,E)\zeta_j^\ell,A(\zeta'-e^X\zeta_j^\ell)\rangle$, whose covariance under the kernel is $\langle L(X,E)\zeta_j^\ell,AL(X,E')\zeta_j^\ell\rangle$, since $A\Qb A=A$; the covariance part of the Gaussian Fisher metric vanishes because $\Qb$ does not depend on $X$. Averaging over $q(\zeta_j^\ell)$ with $\E_q[\zeta_j^\ell\zeta_j^{\ell\top}]=S_j^\ell$ and multiplying by $\beta_{ij}^\ell$ gives~\eqref{dl:eq:fisher}. The second derivative of $\mathcal T_{ij}(e^X)=\tfrac12\beta_{ij}^\ell\E_q\|\zeta_i^{\ell+1}-e^X\zeta_j^\ell\|_A^2$ is the same form minus $\beta_{ij}^\ell\E_q\langle\zeta_i^{\ell+1}-e^X\zeta_j^\ell,AD^2e^X[E,E']\zeta_j^\ell\rangle$, which the Gauss-Newton form omits. For one head the transport term of receiver $i$ is $\sum_j\mathcal T_{ij}(e^{X_{ij}^\ell})$ up to $X$-free terms, so the forms of distinct edges do not couple and $\mathcal I(M)$ is block diagonal. The form $\mathcal I_{ij}^\ell[E,E]=\beta_{ij}^\ell\|A^{1/2}L(X,E)(S_j^\ell)^{1/2}\|_F^2$ vanishes for some $E\ne0$ exactly when $\beta_{ij}^\ell=0$ or $L(X,\cdot)$ has a nonzero kernel, because $A$ and $S_j^\ell$ are positive definite. On the open set where $\mathcal I(M)$ is invertible the field $-\mathcal I(M)^{-1}\nabla_M\mathcal F$ is smooth, so its integral curves exist and are unique; the equation is of first order in $M$, its velocity vanishes exactly where $\nabla_M\mathcal F=0$, and applying $\mathcal I(M)$ gives the force balance. Along a solution $\tfrac{d}{dt}\mathcal F=-\mathcal I(M)[\dot M,\dot M]\le0$, with equality only at a critical point; a solution that reaches a critical point is constant by uniqueness, so $\mathcal F$ decreases strictly along every nonconstant solution, which therefore never returns to an earlier point. Finally $\mathcal I^{-1}H$ is similar to the symmetric matrix $\mathcal I^{-1/2}H\mathcal I^{-1/2}$.
\end{proof}

\begin{proposition}[Non-attainment at semisimple negative eigenvalues]
\label{dl:prop:semisimple}
Let $K\ge2$ and let $f:\R^{K\times K}\to\R$ be differentiable with exactly one stationary point $G^\star$, which is singular. Call an eigenvalue semisimple if all its Jordan blocks have size one. Then no chart link whose negative eigenvalues are all semisimple satisfies $f(G)=\inf_{X\in\gl(K)}f(e^X)$. For $K=2$ and $K=3$ every chart link has this property, so $X\mapsto f(e^X)$ does not attain its infimum over $\gl(K)$.
\end{proposition}

\begin{proof}
Let a chart link $G$ whose negative eigenvalues are all semisimple attain the infimum of $f$ over chart links. It suffices to find a nonempty open set $\mathcal U$ of real matrices $E$ such that, for each $E\in\mathcal U$, $G+tE$ is a chart link for every sufficiently small $|t|$. Indeed, minimality then gives $f(G+tE)\ge f(G)$ for small $t$ of either sign, so $\langle\nabla f(G),E\rangle=0$ for every $E\in\mathcal U$. The linear span of a nonempty open set is $\R^{K\times K}$, so $\nabla f(G)=0$ and $G=G^\star$, which contradicts the invertibility of $G$.

If $G$ has no negative eigenvalue, its spectrum is a compact set at positive distance from $(-\infty,0]$, because $G$ is invertible. By continuity of the roots of the characteristic polynomial, every real matrix near $G$ has no eigenvalue in $(-\infty,0]$ and is a chart link, so $\mathcal U=\R^{K\times K}$ qualifies. Otherwise let $c_1,\dots,c_r$ be the distinct negative eigenvalues of $G$, with multiplicities $n_1,\dots,n_r$. The Jordan blocks of $c_l$ have size one, so their number is $n_l$, which the criterion makes even. The real eigenspaces of the $c_l$ and the sum of the real generalized eigenspaces of the remaining eigenvalues decompose $\R^K$ into invariant subspaces, so that for a real invertible $\Gamma$
\begin{equation*}
\Gamma^{-1}G\Gamma=\operatorname{blockdiag}\big(c_1I_{n_1},\dots,c_rI_{n_r},G_0\big),
\end{equation*}
with $G_0$ real and without eigenvalues in $(-\infty,0]$. For a real $E$ let $\tilde E_{ll}$ be the diagonal block of $\Gamma^{-1}E\Gamma$ at $c_l$, and call $E$ admissible if every $\tilde E_{ll}$ has $n_l$ distinct non-real eigenvalues. The admissible matrices form an open set, because each map $E\mapsto\tilde E_{ll}$ is linear and distinct non-real eigenvalues stay distinct and non-real under small real perturbations. They form a nonempty set, because each $n_l$ is even and the block-diagonal matrix with the $2\times2$ blocks $\big(\begin{smallmatrix}0&-k\\k&0\end{smallmatrix}\big)$, $k=1,\dots,n_l/2$, has the distinct eigenvalues $\pm\mathrm ik$. Let $\mathcal U$ be the set of admissible matrices.

Fix $E\in\mathcal U$ and let $r_0>0$ be smaller than the distance between any two distinct eigenvalues of $G$ and than the distance from the spectrum of $G_0$ to $(-\infty,0]$. Fix $l$, order the blocks so that the one at $c_l$ comes first, let $B$ be the complementary diagonal block of $\Gamma^{-1}G\Gamma$, of which $c_l$ is not an eigenvalue, and split $\Gamma^{-1}E\Gamma$ accordingly into $E_{11}=\tilde E_{ll}$, $E_{12}$, $E_{21}$ and $E_{22}$. For $|\lambda-c_l|\le r_0/2$ and small $|t|$ the matrix $\Delta(t,\lambda)=\lambda I-B-tE_{22}$ is invertible, and the Schur complement gives
\begin{equation*}
\det(\lambda I-G-tE)=\det \Delta(t,\lambda)\det\big((\lambda-c_l)I-tE_{11}-t^2E_{12}\Delta(t,\lambda)^{-1}E_{21}\big).
\end{equation*}
With $\lambda=c_l+t\nu$ the second factor is $t^{n_l}\phi(t,\nu)$, where $\phi(t,\nu)=\det\big(\nu I-E_{11}-tE_{12}\Delta(t,c_l+t\nu)^{-1}E_{21}\big)$ is analytic in $(t,\nu)$ near $t=0$, and $\phi(0,\cdot)$ is the characteristic polynomial of $E_{11}$, whose roots $\nu_1,\dots,\nu_{n_l}$ are simple and non-real. The implicit function theorem at these simple roots gives continuous $\nu_k(t)$ for small $|t|$ with $\phi(t,\nu_k(t))=0$ and $\nu_k(0)=\nu_k$, and these stay distinct and non-real. For small $t\ne0$ the numbers $c_l+t\nu_k(t)$ are therefore $n_l$ distinct eigenvalues of $G+tE$ within $r_0/2$ of $c_l$, and they are non-real, with imaginary parts $t\operatorname{Im}\nu_k(t)\ne0$. By continuity of the roots of the characteristic polynomial, for small $|t|$ exactly $n_l$ eigenvalues of $G+tE$, counted with multiplicity, lie within $r_0/2$ of $c_l$, so these are all of them. The same continuity places every eigenvalue of $G+tE$ within $r_0/2$ of the spectrum of $G$, and those not near any $c_l$ lie near the spectrum of $G_0$, off $(-\infty,0]$. Hence, for every small $t\ne0$ of either sign, $G+tE$ has no eigenvalue in $(-\infty,0]$ and is a chart link by the criterion, and for $t=0$ it is $G$.

For $K\le3$, a negative eigenvalue of a chart link has Jordan blocks of each size occurring an even number of times, so its multiplicity is at least two. Multiplicity two with one block of size two would make that block occur once, and multiplicity three cannot be split into blocks of each size occurring an even number of times. The negative eigenvalue therefore has multiplicity two and two blocks of size one, so it is semisimple. Finally, $f(e^X)$ attains the infimum over $\gl(K)$ at $X$ exactly when $e^X$ attains the infimum of $f$ over chart links.
\end{proof}

For the population statement, fix an edge $(i\leftarrow j)$ at layer $\ell$ with one head, write $a=(i,\ell+1)$ and $b=(j,\ell)$ for its vertices, and let a population of contexts $\omega$ carry mean-field recognition laws $q_\omega$ and routing weights $\beta_\omega=\beta_{ij,\omega}^\ell>0$, with $0<\E_\omega\beta_\omega<\infty$. The routing-weighted population law $\mathbb P$ of $(\zeta_a,\zeta_b)$ is defined by $\E_{\mathbb P}[g]=\E_\omega[\beta_\omega\E_{q_\omega}g(\zeta_a,\zeta_b)]/\E_\omega[\beta_\omega]$, and when the weights do not vary across the population it is the mixture $\E_\omega[q_\omega]$, a common law of the states at all vertices.

\begin{proposition}[Population regression links and their curvature]
\label{dl:prop:regression}
\begin{enumerate}
\item[(i)] Assume that $\mathbb P$ has finite second moments and that $M_{bb}=\E_{\mathbb P}[\zeta_b\zeta_b^\top]$ is invertible, and write $M_{ab}=\E_{\mathbb P}[\zeta_a\zeta_b^\top]$. The population average $\E_\omega[\mathcal T_{ij,\omega}(G)]=\tfrac12\E_\omega[\beta_\omega]\E_{\mathbb P}\|\zeta_a-G\zeta_b\|_A^2$ of the transport term is a strictly convex quadratic function of $G\in\R^{K\times K}$ whose unique minimizer $M_{ab}M_{bb}^{-1}$ solves the population average $GM_{bb}=M_{ab}$ of the normal equations $GS_{j,\omega}^\ell=\xi_{i,\omega}^{\ell+1}\xi_{j,\omega}^{\ell\top}$ of the single contexts. If $\mathbb P$ is centered, the minimizer is the population regression link $W_{ab}=C_{ab}C_{bb}^{-1}$, with $C_{xy}=\operatorname{Cov}_{\mathbb P}(\zeta_x,\zeta_y)$.
\item[(ii)] Let $\mathbb P$ be a centered law of the states at vertices $a$, $b$ and $c$ with finite second moments and invertible $C_{bb}$ and $C_{cc}$, let $W_{xy}=C_{xy}C_{yy}^{-1}$, and put $C_{ac\mid b}=C_{ac}-C_{ab}C_{bb}^{-1}C_{bc}$. Then
\begin{equation}
W_{ab}W_{bc}=W_{ac}-C_{ac\mid b}C_{cc}^{-1},
\label{dl:eq:partialcov}
\end{equation}
so the route through $b$ composes to the direct regression link exactly when the partial cross-covariance $C_{ac\mid b}$ vanishes; for a jointly Gaussian $\mathbb P$ this is the conditional independence of $\zeta_a$ and $\zeta_c$ given $\zeta_b$. Consequently, for the regression links of one centered law, the triangle holonomy $F_{ijk}=W_{ik}^{-1}W_{ij}W_{jk}$ of a layer-independent field with $W_{ik}$ invertible is $I-W_{ik}^{-1}C_{ik\mid j}C_{kk}^{-1}$, which is trivial exactly when $C_{ik\mid j}=0$, and the loop~\eqref{dl:eq:loops} built from regression links, where its two routes are invertible, is trivial exactly when the partial cross-covariances of $(i,\ell+2)$ and $(k,\ell)$ given $(j,\ell+1)$ and given $(k,\ell+1)$ coincide.
\end{enumerate}
\end{proposition}

\begin{proof}
(i) By the definition of $\mathbb P$, $\E_\omega[\beta_\omega\E_{q_\omega}\|\zeta_a-G\zeta_b\|_A^2]=\E_\omega[\beta_\omega]\E_{\mathbb P}\|\zeta_a-G\zeta_b\|_A^2$, and
\begin{equation*}
\E_{\mathbb P}\|\zeta_a-G\zeta_b\|_A^2=\tr(AM_{aa})-2\tr(AGM_{ab}^\top)+\tr(AGM_{bb}G^\top),
\end{equation*}
whose gradient in $G$ is $2A(GM_{bb}-M_{ab})$ and whose Hessian $E\mapsto2AEM_{bb}$ is positive definite, as in the proof of Proposition~\ref{dl:prop:linear}. The unique stationary point is $M_{ab}M_{bb}^{-1}$. Under the mean-field law $q_\omega$, $\E_{q_\omega}[\zeta_a\zeta_b^\top]=\xi_{i,\omega}^{\ell+1}\xi_{j,\omega}^{\ell\top}$ and $\E_{q_\omega}[\zeta_b\zeta_b^\top]=S_{j,\omega}^\ell$, so $GM_{bb}=M_{ab}$ is the routing-weighted average of the single-context normal equations. For a centered law the second moments are the covariances.

(ii) Multiplication gives $W_{ab}W_{bc}=C_{ab}C_{bb}^{-1}C_{bc}C_{cc}^{-1}=(C_{ac}-C_{ac\mid b})C_{cc}^{-1}$, which is~\eqref{dl:eq:partialcov}, and since $C_{cc}$ is invertible the composite equals $W_{ac}$ exactly when $C_{ac\mid b}=0$. For a jointly Gaussian law with $C_{bb}$ invertible, the conditional law of $(\zeta_a,\zeta_c)$ given $\zeta_b$ is Gaussian with a covariance that does not depend on the value of $\zeta_b$, namely the Schur complement of $C_{bb}$, whose off-diagonal block is $C_{ac\mid b}$; jointly Gaussian vectors are independent exactly when their cross-covariance vanishes \cite{dl:lauritzen1996}. For the triangle, multiply~\eqref{dl:eq:partialcov} with $(a,b,c)=(i,j,k)$ by $W_{ik}^{-1}$ on the left. For the loop, put $a=(i,\ell+2)$ and $c=(k,\ell)$. By~\eqref{dl:eq:partialcov} the route through $b'=(j,\ell+1)$ composes to $(C_{ac}-C_{ac\mid b'})C_{cc}^{-1}$ and the route through $b''=(k,\ell+1)$ to $(C_{ac}-C_{ac\mid b''})C_{cc}^{-1}$, and $W_{ijk}^\ell$, the inverse of the second times the first, is the identity exactly when the two routes agree, that is, when $C_{ac\mid b'}=C_{ac\mid b''}$.
\end{proof}

The partial cross-covariance records conditional linear dependence beyond the middle vertex; for non-Gaussian laws its vanishing does not by itself imply conditional independence. The regression links need cross-covariances between distinct vertices, which vanish under a single mean-field recognition law and appear only in the population law or in a correlated recognition family.

\begin{proposition}[Monotone maximum a posteriori path of a scalar link coordinate]
\label{dl:prop:kink}
Let the link coordinate be scalar, $x^\ell\in\R$ for $\ell=0,\dots,L$ with $L\ge1$, hold the ends at $x^0=a<b=x^L$, let $m>0$, and let $V:\R\to\R$ be continuous with $V\ge V(a)$ on $(-\infty,a]$ and $V\ge V(b)$ on $[b,\infty)$, as for a double well with minima at $a$ and $b$. Then
\begin{equation*}
\mathcal E(x)=\frac m2\sum_{\ell=0}^{L-1}\big(x^{\ell+1}-x^\ell\big)^2+\sum_{\ell=1}^{L-1}V(x^\ell),
\end{equation*}
the objective of Proposition~\ref{dl:prop:inertia} with fixed ends, attains its minimum over the interior values $x^1,\dots,x^{L-1}$, and every minimizing path is nondecreasing in $\ell$. If moreover $V$ is differentiable with $V'(a)=V'(b)=0$, every minimizing path is strictly increasing.
\end{proposition}

\begin{proof}
Replacing every interior value by its clamp to $[a,b]$ does not increase $\mathcal E$: the clamp is $1$-Lipschitz, so no increment grows, and by the hypothesis on $V$ outside $[a,b]$ no potential term grows. If some interior value lies outside $[a,b]$, take a maximal run of consecutive such values. The ends lie in $[a,b]$, so the run is entered by a step from a value $y\in[a,b]$, and clamping strictly shortens that step; the decrease is then strict. The continuous function $\mathcal E$ attains its minimum on the compact box $[a,b]^{L-1}$, this minimum is global, and every global minimizer takes its values in $[a,b]$.

Let $z_0\le z_1\le\dots\le z_L$ be the values of a minimizer, ends included, in increasing order, so that $z_0=a$, $z_L=b$, and $z_1,\dots,z_{L-1}$ are the interior values as a multiset; let $g_k=z_{k+1}-z_k$. No value lies strictly inside an interval $(z_k,z_{k+1})$, so each step from $x^\ell$ to $x^{\ell+1}$ covers a union of whole gaps, and its square is at least the sum of the squares of the gaps it covers. A path from $a$ to $b$ crosses each gap of positive length an odd number of times, so $\sum_\ell(x^{\ell+1}-x^\ell)^2\ge\sum_kg_k^2$, which is the kinetic sum of the increasing rearrangement $(z_0,\dots,z_L)$. A decreasing step crosses some gap of positive length downward, and that gap is then crossed at least three times, so the inequality is strict. The rearrangement keeps the ends and the multiset of interior values, and hence the potential sum. A path with a decreasing step therefore has a strictly better competitor, and every minimizer is nondecreasing.

For strictness, let $V$ be differentiable with $V'(a)=V'(b)=0$. The interior values of a minimizer are free variables, so Eq.~\eqref{dl:eq:inertia} holds at every $0<\ell<L$. Suppose a minimizer has a maximal run $x^p=\dots=x^q=x$ with $q\ge p+1$. If $0<p$ and $q<L$, Eq.~\eqref{dl:eq:inertia} at $p$ and at $q$ gives $m(x-x^{p-1})=-V'(x)=m(x-x^{q+1})$, whereas monotonicity and maximality give $x^{p-1}<x<x^{q+1}$, a contradiction. If $p=0$, then $x=a$ and $q<L$ because $a<b$, and Eq.~\eqref{dl:eq:inertia} at the interior index $q$ gives $m(a-x^{q+1})=-V'(a)=0$, against maximality. The case $q=L$ is symmetric.
\end{proof}

\section{Numerical checks}
\label{dl:app:checks}

The checks run in float64 with NumPy and SciPy. Check~1 evaluates the eigenvalue density and the spectral curvature on $I+E_{12}$ (both zero) and the Frobenius defect ($0.5$), and shows conjugates of the shear converging to $I$. Check~2 confirms $\Phi(Q)=K-\tr Q$ on $200$ random orthogonal matrices to $1.8\times10^{-15}$ and the small-field expansion~\eqref{dl:eq:delta} with relative error $10^{-3}$, $3.4\times10^{-4}$ and $3.8\times10^{-5}$ at scales $10^{-1}$, $10^{-2}$ and $10^{-3}$. Check~3 compares the transport term~\eqref{dl:eq:transport} with a Monte Carlo estimate over $4\times10^5$ joint samples for two heads and three sources; the relative difference, $9.5\times10^{-5}$, is at the level of the sampling error, and the sampled mean and covariance of $m$ match~\eqref{dl:eq:moments} and~\eqref{dl:eq:V} to $1.3\times10^{-3}$ and $1.4\times10^{-2}$. Check~4 confirms the one-head identity~\eqref{dl:eq:transport_onehead} exactly. Check~5 compares the gradient of Proposition~\ref{dl:prop:field} with central differences ($3.0\times10^{-10}$) and confirms that the unregularized minimizer over all matrices has rank one. Check~6 compares the linearized links~\eqref{dl:eq:linear} with the exact stationary links of a four-token transition with triangle loops; the relative difference is $2.7\times10^{-2}$, $8.5\times10^{-3}$ and $2.9\times10^{-3}$ at source scales $10^{-1}$, $3\times10^{-2}$ and $10^{-2}$, and the operator's smallest eigenvalue is $2.014$ against the bound $\sigma^{-2}=2$. Check~7 confirms the electric-loop expansion with relative error decreasing linearly in the scale. Check~8 confirms orthogonal gauge invariance of the induced action to machine precision, the hopping series to $10^{-15}$, the leading Wilson coefficient through the ratios quoted in Section~\ref{dl:sec:induced}, the absence of an induced action for orthogonal links on a tree, its presence for a symmetric generator ($0.237$), its response to a unipotent link ($0.045$), and the lower bound over $500$ random large $\GL(2)$ configurations. Check~9, for a scalar path with $L=40$, $m=25$, ends $\mp1$ and $V(x)=(x^2-1)^2/4$, confirms~\eqref{dl:eq:inertia} at the computed minimizer to $7\times10^{-7}$ and finds the computed path strictly increasing, an instance of Proposition~\ref{dl:prop:kink}. Check~10 confirms the static limit of Remark~\ref{dl:rem:static} to $5.6\times10^{-16}$.

\begin{thebibliography}{99}
\bibitem{dl:wilson1974} K. G. Wilson, Confinement of quarks, \emph{Physical Review D} 10, 2445--2459, 1974.
\bibitem{dl:kogut1975} J. Kogut and L. Susskind, Hamiltonian formulation of Wilson's lattice gauge theories, \emph{Physical Review D} 11, 395--408, 1975.
\bibitem{dl:creutz1983} M. Creutz, \emph{Quarks, Gluons and Lattices}, Cambridge University Press, 1983.
\bibitem{dl:montvay1994} I. Montvay and G. M\"unster, \emph{Quantum Fields on a Lattice}, Cambridge University Press, 1994.
\bibitem{dl:kazakov1993} V. A. Kazakov and A. A. Migdal, Induced QCD at large $N$, \emph{Nuclear Physics B} 397, 214--238, 1993.
\bibitem{dl:lauritzen1996} S. L. Lauritzen, \emph{Graphical Models}, Oxford University Press, 1996.
\bibitem{dl:culver1966} W. J. Culver, On the existence and uniqueness of the real logarithm of a matrix, \emph{Proceedings of the American Mathematical Society} 17(5), 1146--1151, 1966.
\bibitem{dl:higham2008} N. J. Higham, \emph{Functions of Matrices: Theory and Computation}, SIAM, 2008.
\bibitem{dl:khan2023} M. E. Khan and H. Rue, The Bayesian learning rule, \emph{Journal of Machine Learning Research} 24(281), 1--46, 2023.
\bibitem{dl:onsager1953} L. Onsager and S. Machlup, Fluctuations and irreversible processes, \emph{Physical Review} 91, 1505--1512, 1953.
\bibitem{dl:friston2008dem} K. J. Friston, N. Trujillo-Barreto and J. Daunizeau, DEM: A variational treatment of dynamic systems, \emph{NeuroImage} 41, 849--885, 2008.
\bibitem{dl:schmidhuber1992} J. Schmidhuber, Learning to control fast-weight memories: An alternative to dynamic recurrent networks, \emph{Neural Computation} 4, 131--139, 1992.
\end{thebibliography}

\end{document}
